\documentclass[aspectratio=169,11pt]{beamer}

% Japanese support (compile with lualatex)
\usepackage{luatexja}
\usepackage{amsmath,amssymb}
\usepackage{bm}

\usetheme{Madrid}
\usecolortheme{default}

\title{Beamer Sample}
\subtitle{Diffusion Equation}
\author{T. Toida}
\institute{Simulation Engineering}
\date{\today}

\begin{document}

\begin{frame}
  \titlepage
\end{frame}

\begin{frame}{Agenda}
  \tableofcontents
\end{frame}

\section{Model}

\begin{frame}{Diffusion Equation}
  Consider the 1D diffusion equation:
  \[
    \frac{\partial u}{\partial t} = K\frac{\partial^2 u}{\partial x^2},
    \quad 0 < x < 1,\; t > 0,\; K > 0.
  \]
  Example boundary conditions (Dirichlet):
  \[
    u(0,t)=u(1,t)=0.
  \]
\end{frame}

\section{Fourier Expansion}

\begin{frame}{Fourier Sine Expansion}
  Under homogeneous Dirichlet BC, we use
  \[
    u(x,t)=\sum_{k=1}^{\infty}\hat{u}(k,t)\sin(k\pi x).
  \]
  Coefficients are obtained by
  \[
    \hat{u}(k,t)=2\int_0^1 u(x,t)\sin(k\pi x)\,dx.
  \]
\end{frame}

\begin{frame}{Mode Decay}
  Each Fourier mode satisfies
  \[
    \frac{d\hat{u}(k,t)}{dt}=-(k\pi)^2K\,\hat{u}(k,t),
  \]
  therefore
  \[
    \hat{u}(k,t)=\hat{u}(k,0)e^{-(k\pi)^2Kt}.
  \]
  \begin{itemize}
    \item Large $k$ (fine structures) decay faster.
    \item The solution becomes smoother as time evolves.
  \end{itemize}
\end{frame}

\section{Summary}

\begin{frame}{Summary}
  \begin{enumerate}
    \item Diffusion smooths spatial variations.
    \item Fourier decomposition gives mode-wise dynamics.
    \item Beamer is suitable for math-heavy slides.
  \end{enumerate}
\end{frame}

\begin{frame}
  \centering
  {\Huge Thank you!}
\end{frame}

\end{document}
